\documentclass[12pt,a4paper]{exam}

% --------------------
% Packages
% --------------------
\usepackage[utf8]{inputenc}
\usepackage[T1]{fontenc}
\usepackage[french]{babel}
\usepackage{amsmath,amssymb}
\usepackage{geometry}
\usepackage{tikz}
\usetikzlibrary{angles,quotes}
\usetikzlibrary{calc}
\usepackage{pgfplots}
\pgfplotsset{compat=1.18}
\geometry{margin=2cm}

% Pour les coordonnées de vecteurs :
\newcommand*{\Coord}[2]{%
  \ensuremath{
    \begin{pmatrix}
      #1 \\[0pt]
      #2
    \end{pmatrix}}
}

% --------------------
% En-tête / pied (exam)
% --------------------
\header{\textbf{NOM Prénom : }}{}{2\up{nde}}
\footer{}{\thepage}{}

% --------------------
\begin{document}

\begin{center}
    \Large\textbf{Devoir Surveillé\\Informations chiffrées et calculs}
\end{center}

\vspace{0.3cm}

\noindent
\textbf{Durée :} 1h \hfill
\textbf{Calculatrice :} autorisée

\vspace{0.5cm}

% ========================================================
%             Exo : graphique
% ========================================================
\noindent
\textbf{Exercice 1 (\ \ \ \  / 2 points )}

\vspace{10pt}
Dans un lycée, 400 élèves sont en Seconde. 30\% des élèves s'orientent en filière 
technologique en fin de 
Seconde.

Parmi eux, 40\% vont en série STMG.

\begin{questions}

\question[2] Déterminer combien d'élèves vont en STMG.

\end{questions}

\vspace{1cm}

% ========================================================
%             Exo : graphique
% ========================================================
\noindent
\textbf{Exercice 2 (\ \ \ \  / 3 points )}

\vspace{10pt}
Le prix du litre d'essence est passé de 1,67 euros en janvier 2022 à 2,15 euros en juin 2022.

\begin{questions}

\question[1] Calculer le taux d'évolution d'un tel changement. Arrondir à $10^-2$ près.

\question[2] Quelle différence de prix cette variation représente-t-elle sur un plein de 30 litres ?

\end{questions}

\vspace{1cm}

% ========================================================
%             Exo : graphique
% ========================================================
\noindent
\textbf{Exercice 3 (\ \ \ \  / 3 points )}

\vspace{10pt}
Lors de la première démarque des soldes, un magasin annonce -30\% sur toutes les chemises. 

Lors de la deuxième démarque, il annonce -10\% supplémentaires sur tous les produits soldés.

\begin{questions}

\question[2] Quel est le taux d'évolution global en pourcentage du prix d'une chemise ?

\question[1] Avant la première démarque, le prix d'une chemise était de 60€. Quel est alors son prix final ?

\end{questions}

\vspace{1cm}

% ========================================================
%             Exo : graphique
% ========================================================
\noindent
\textbf{Exercice 4 (\ \ \ \  / 2 points )}

\vspace{10pt}
Dans une réserve naturelle, la population de tigres diminue de 18\%. 

\begin{questions}

\question[2] De quel pourcentage faut-il qu'elle augmente pour revenir à sa valeur initiale ?
Arrondir à 0,01\% près.

\end{questions}

\vspace{1cm}

% ========================================================
%             Exercice 2
% ========================================================
\noindent
\textbf{Exercice 5 : Du calcul ( \ \ \ / 6 points )}

\vspace{10pt}
\textbf{\underline{Les prochaines question sont à faire sans calculatrice, en détaillant vos calculs.}}

\vspace{5pt}
\textbf{\underline{ATTENTION :}} si vous ne détaillez pas vos étapes de calculs, je saurai que vous avez 
utilisé la calculatrice.
  
\vspace{10pt}
\begin{questions}

\question[4] Mettre sous forme simplifiée les expressions suivantes :

\vspace{5pt}
\hspace{1.85cm} a) $\dfrac{13}{30} - \dfrac{7}{15} + \dfrac{5}{3}$ \hspace{2cm} b) $\dfrac{2}{11} + 2$

\vspace{10pt}
\hspace{1.85cm} c) $\dfrac{2}{3}\left(2+\dfrac{3}{4}\right)$ \hspace{2.3cm} 
d) $\dfrac{7}{4} \div 2 - \dfrac{1}{6} \times \dfrac{2}{-3}$

\vspace{10pt}
\question[1] Écrire sous la forme $a\sqrt{b}$ où $a$ et $b$ sont des nombres entiers positifs, $b$ étant le plus 
\textbf{petit possible}.

\vspace{5pt}
\hspace{1.85cm} a) $\sqrt{50}$ \hspace{2cm} b) $\sqrt{54}$

\vspace{10pt}
\question[1] Écrire les nombres suivants en notation scientifique :

\vspace{5pt}
\hspace{1.85cm} a) $232$ \hspace{2cm} b) $0,958$

\vspace{10pt}
\question[2] Écrire sous la forme $a^n$, où $a$ est un nombre relatif et $n$ un entier relatif :

\vspace{5pt}
\hspace{1.85cm} a) $\dfrac{3^8}{3^{-4}}$ \hspace{2cm} b) $\dfrac{(-6)^4}{3^4}$

\vspace{5pt}
\hspace{1.85cm} c) $8 \times 8^{-3}$ \hspace{1.4cm} d) $((-3,7)^{-2})^4$

\vspace{10pt}
\question[1] Montrer que pour tout entier naturel $n \neq 0$,

\vspace{5pt}
\hspace{4cm} $\dfrac{1}{n+1} - \dfrac{1}{n} = \dfrac{-1}{n(n+1)}$

\vspace{10pt}
\question[1] Écrire sous la forme $a + b\sqrt{c}$ où $c$ est le plus \textbf{petit possible} :

\vspace{5pt}
\hspace{1.85cm} a) $5\sqrt{7}(-2 + 4\sqrt{7})$ \hspace{2cm} b) $(-9\sqrt{11} + 12)(4 - 3\sqrt{11})$

\vspace{40pt}
\begin{center}
    \textit{-- Fin du devoir --}
\end{center}

\vspace{40pt}
\question[2] \textbf{\underline{BONUS}} Soit $E = \dfrac{5}{\sqrt{2}+\sqrt{18}} + \dfrac{3}{\sqrt{2} - \sqrt{18}}$.

\vspace{5pt}
Écrire le nombre $E$ sous la forme $a\sqrt{b}$ où $a$ est une fraction irréductible et $b$ est un nombre entier.

\end{questions}


\end{document}
